\documentclass[11pt,letterpaper]{article}

\usepackage[spanish,es-noshorthands,es-tabla]{babel}
\usepackage{fontspec}
\setmainfont{Source Serif Pro}
\usepackage{microtype}
\usepackage{geometry}
\usepackage{enumitem}
\usepackage{xcolor}
\usepackage{csquotes}
\usepackage{float}
\usepackage{booktabs}
\usepackage{tikz}
\usetikzlibrary{arrows.meta,positioning,fit,backgrounds}
\usepackage{xurl}
\usepackage[hidelinks]{hyperref}
\usepackage[backend=biber,style=apa,sorting=nyt,maxcitenames=2]{biblatex}
\DeclareLanguageMapping{spanish}{spanish-apa}

\geometry{margin=2.3cm}
\hbadness=10000
\setlength{\emergencystretch}{1em}
\setlength{\parindent}{0pt}
\setlength{\parskip}{0.5em}
\setlist{nosep}
\addbibresource{referencias.bib}
\AtBeginBibliography{\sloppy\small}
\setlength{\bibitemsep}{0.3em}

% Colores: terracota para las etapas del círculo vicioso, azul oscuro para
% rótulos y flechas, gris para las notas con los datos de respaldo.
\definecolor{cCiclo}{RGB}{186,86,30}
\definecolor{cEje}{RGB}{26,58,102}
\definecolor{cGris}{RGB}{120,116,110}

\tikzset{
  % los rótulos de la figura no se dividen con guion
  every node/.append style={execute at begin node={\hyphenpenalty=10000}},
  etapa/.style={draw=cCiclo, fill=cCiclo!9, thick, rounded corners=3pt,
    text width=3.6cm, minimum height=2.3cm, align=center, inner sep=5pt,
    font=\small},
  dato/.style={font=\scriptsize, text=cGris, text width=3.6cm, align=center},
  flecha/.style={cEje, very thick, -{Stealth[length=3mm]}},
}

\begin{document}

% --- Portada -----------------------------------------------------------------
\begin{titlepage}
\centering
\vspace*{3.2cm}
{\large\scshape Gestión del Desarrollo Tecnológico:\\
Estrategias y Análisis de Portfolio\par}
\vspace{0.5cm}
{\normalsize Módulo 7. Foro 7.4: Innovación abierta e innovación\\
en el sistema militar de salud\par}
\vspace{2.6cm}
{\color{cEje}\rule{0.55\textwidth}{0.8pt}\par}
\vspace{0.7cm}
{\LARGE\bfseries Cuando el cliente es el Estado\par}
\vspace{0.5cm}
{\large Por qué la burocracia aleja a los mejores innovadores\\
y encarece lo que compramos todos\par}
\vspace{0.7cm}
{\color{cEje}\rule{0.55\textwidth}{0.8pt}\par}
\vfill
{\large Carlos Luis Rojas Aragonés\par}
\vspace{0.3cm}
{\normalsize Chief Technology Officer Program\par}
\vspace{0.3cm}
{\normalsize 21 de octubre de 2026\par}
\vspace*{1.5cm}
\end{titlepage}

% --- Cuerpo ------------------------------------------------------------------
Trabajar con el gobierno es difícil en términos generales, y cuando el tema
es de seguridad lo es todavía más. El artículo de MIT Sloan lo ilustra con
casos concretos que recoge Tina Srivastava, exingeniera jefe de programas de
guerra electrónica en Raytheon: un contrato de la Fuerza Aérea con una
startup, para necesidades \enquote{urgentes y críticas} de seguridad
nacional, tardó tres años en arrancar, y el oficial de contratación dejó de
atender a la empresa hasta que contrató a un abogado
\parencite{somers2019}.

\section*{1. Dos consecuencias que se alimentan entre sí}

En mi opinión, esa fricción produce dos efectos:

\begin{enumerate}[label=\arabic*., leftmargin=1.6em, itemsep=0.35em]
  \item \textbf{Las empresas simplemente evitan participar} por el exceso de
    burocracia y lo lento que se mueve todo. El dato lo confirma: las
    pequeñas empresas con contratos del Departamento de Defensa pasaron de
    42\,723 en 2011 a 24\,296 en 2020, un 43\,\% menos. La GAO recogió la
    razón de boca de los propios empresarios: muchas prefieren el sector
    privado porque trabajar con Defensa implica retrasos impredecibles en la
    adjudicación, las autorizaciones de seguridad, el inicio de los trabajos
    y los pagos \parencite{gao2021}. Y eso suponiendo que el proceso sea
    transparente y que todos compitan en igualdad de condiciones, algo que
    no siempre se cumple: el artículo cuenta cómo se excluye a quien no
    tiene \enquote{historial de desempeño} y cómo Tenebraex gastó cientos
    de miles de dólares en abogados para cobrar por su propia patente,
    después de que el Ejército adjudicara el contrato solo a su socio, una
    filial de Raytheon \parencite{somers2019}.
  \item \textbf{Las pocas empresas que sí participan inflan el precio},
    porque saben lo que implica ganar un proyecto con el gobierno y porque
    tienen menos competencia. La industria de defensa estadounidense pasó de
    51 a 5 contratistas principales en la década de 1990 \parencite{dod2022}.
    Con proveedores únicos, el Inspector General encontró que TransDigm
    obtuvo ganancias excesivas en 46 de 47 repuestos, con márgenes de hasta
    4\,451\,\% \parencite{dodig2019}, y en 2024 que la Fuerza Aérea aceptó un
    sobreprecio del 7\,943\,\% en dispensadores de jabón para el C-17
    \parencite{dodig2024}. Las empresas cuestionan la metodología, pero el
    patrón se repite.
\end{enumerate}

Esos dos factores se refuerzan (Figura~\ref{fig:ciclo}): hay gente muy
buena que no quiere o no puede participar, y todos terminamos pagando el
costo asociado a cada proyecto en forma de sobreprecios. Peor aún, cada
abuso descubierto justifica más controles, que vuelven el proceso más lento
y espantan a otros participantes nuevos.

\begin{figure}[H]
\centering
\begin{tikzpicture}[node distance=1.5cm and 2.4cm]
  \node[etapa] (buro) {\textbf{Burocracia y lentitud}\\registros, cláusulas,
    secreto, plazos de años};
  \node[etapa, right=of buro] (pocos) {\textbf{Pocos participantes}\\las
    startups y pymes prefieren el sector privado};
  \node[etapa, below=of pocos] (precio) {\textbf{Sobreprecios}\\proveedores
    únicos fijan el precio};
  \node[etapa, below=of buro] (control) {\textbf{Más controles}\\cada abuso
    se responde con nuevas reglas};

  \draw[flecha] (buro) -- (pocos);
  \draw[flecha] (pocos) -- (precio);
  \draw[flecha] (precio) -- (control);
  \draw[flecha] (control) -- (buro);

  \node[dato, above=0.15cm of buro] {Un contrato tardó 3 años en arrancar};
  \node[dato, above=0.15cm of pocos] {Pymes en Defensa: \(-\)43\,\% (2011 a 2020)};
  \node[dato, below=0.15cm of precio] {De 51 a 5 contratistas principales;
    márgenes de hasta 4\,451\,\%};
  \node[dato, below=0.15cm of control] {Interpretación propia};
\end{tikzpicture}
\caption{El círculo vicioso de la innovación con el gobierno: la fricción
reduce la competencia, la falta de competencia encarece las compras y los
abusos traen más fricción. Elaboración propia con datos de
\textcite{somers2019}, \textcite{gao2021}, \textcite{dod2022} y
\textcite{dodig2019}.}
\label{fig:ciclo}
\end{figure}

\section*{2. El sistema militar de salud no puede innovar solo}

El enlace al artículo de Deloitte está roto; lo leí en una copia de
Internet Archive \parencite{deloitte2020}. Propone diez innovaciones:
telemedicina, realidad aumentada, impresión 3D, cirugía robótica, atención
centrada en el paciente, dispositivos vestibles, inteligencia aumentada,
blockchain, medicina de precisión y medicina regenerativa. Lo que me llama la
atención es que ninguna nació en el ámbito militar: todas vienen del mercado
civil. El propio artículo lo admite al decir que los pacientes militares
tienen las mismas expectativas que los de cualquier otro sistema de salud, y
recomienda buscar conocimiento en \enquote{fuentes no tradicionales},
construir un ecosistema de socios y empezar con pilotos pequeños
\parencite{deloitte2020}.

Es decir, la receta de Deloitte \emph{es} innovación abierta, y por eso
choca de frente con las barreras que describe MIT Sloan. Una startup de
telemedicina o de dispositivos vestibles que puede venderle a cualquier
hospital privado en semanas difícilmente aceptará esperar años y firmar
cientos de cláusulas para venderle al Departamento de Defensa.

\section*{3. Tecnología comercial con fines militares}

La guerra en Ucrania es el mejor ejemplo de lo que pasa cuando la tecnología
comercial entra sin el filtro tradicional. Starlink, un servicio de internet
satelital para consumidores, se convirtió en la red de comunicaciones del
frente, y en 2023 el Pentágono terminó comprando terminales adicionales para
Ucrania \parencite{techcrunch2023}. Los drones comerciales pasaron a ser
munición desechable: en 2023 Ucrania perdía unos 10\,000 al mes
\parencite{watling2023}, y en 2024 su Ministerio de Defensa compró más de
1,5 millones, el 96\,\% de fabricantes nacionales \parencite{euromaidan2025}.
La ventaja no vino de una tecnología superior, sino de iterar más rápido y
más barato que el adversario.

El caso contrario es IVAS, el visor de combate basado en el HoloLens de
Microsoft, con un contrato de hasta 21\,880 millones de dólares. En 2025 el
contrato pasó a Anduril y, según la GAO, ninguna de las 10\,000 unidades
desarrolladas llegará a usarse en campo \parencite{crs2026}. Adaptar un
producto comercial no basta si se hace con el ciclo de adquisición
tradicional, que hoy promedia más de 12 años para entregar una capacidad
\parencite{gao2026}.

Hay intentos de romper el círculo. La Defense Innovation Unit (DIU) usa
contratos de prototipo con menos requisitos y aspira a adjudicar en 60 a 90
días; en la práctica tardó 142 en 2022 \parencite{diu2022}, lejos de su
meta, pero muy lejos también de los años habituales. El resultado se nota en
quién participa: el 87\,\% de sus proveedores son no tradicionales y el
38\,\% le vendía a Defensa por primera vez \parencite{diu2024}. Cuando se
baja la barrera, aparecen los innovadores que antes no querían o no podían
entrar.

\section*{4. Lo que me llevo como CTO}

El problema del gobierno no es la falta de ideas, sino que su proceso de
compra selecciona a quienes saben navegar la burocracia en lugar de a
quienes mejor resuelven el problema. Lo mismo ocurre, a menor escala, en
empresas grandes con procesos de compra pesados: si homologar a un proveedor
nuevo toma seis meses, terminamos trabajando siempre con los mismos y
pagando lo que nos piden. La lección es diseñar la puerta de entrada (pilotos
pequeños, contratos simples, pagos rápidos) con el mismo cuidado que el
proyecto, porque de ella depende quién se anima a tocarla.

\printbibliography[title={Referencias}]

\end{document}
